\begin{abstract}
Electricity load shedding remains a critical operational challenge in developing power system infrastructures. This paper presents SmartGrid Sentinel, a comparative machine learning and deep sequence forecasting framework designed to predict grid instability and load shedding risks at a localized Upazila level. Utilizing a dataset containing 65,983 records spanning 41 Upazilas in the Sylhet and Chattogram divisions of Bangladesh, we formulate a 2-hour ahead sequence forecasting task. We evaluate twelve architectures, ranging from classical machine learning baselines to recurrent networks and state-of-the-art time-series Transformers. To prevent lookahead bias and target leakage, we implement a strict chronological train-test split (80/20) and restrict feature scaling parameter fitting to the training set only. The PyTorch-based Informer model achieves a Weighted F1 score of 88.44\% and an Accuracy of 87.41\%. Although standard GRU achieves a raw accuracy of 87.69\%, it suffers from severe validation loss divergence (overfitting) after 20 training epochs. The Informer model demonstrates superior validation stability, handles class imbalance robustly—yielding a recall of 79.00\% on the minority Medium-Risk class—and scales efficiently. Finally, we integrate a Rule-Based Explainability Layer to translate high-risk telemetry predictions into human-readable action protocols for grid operators and general citizens.
\end{abstract}

\begin{IEEEkeywords}
Smart Grid, Load Shedding Prediction, Time Series Forecasting, Informer, Transformer, Bangladesh.
\end{IEEEkeywords}
